\section{参考文献：薄膜光学与不确定性方法}
\label{sec:references}

前章提出光学标定与重复测量的方向，本节说明界面反射和区间检验的方法来源，区分物理推导与观测核对。空气、外延层和衬底构成三层介质，两种偏振的界面导纳及往返场在模型建立部分由边界条件和逐束叠加推出。表\ref{tab:available-info}与表\ref{tab:03}区分了观测给定量和待估参数；附件没有折射率曲线，通用光学关系中的材料参数保留为待估量。
% src:交接/数据档案.json:未随附件提供的建模输入

连续留段把相邻谱点保留在同一区段，再检查未参与拟合的光谱。问题二和问题三分别使用$320$个和$640$个测试点核对反射率预测区间，表\ref{tab:09}汇总其经验覆盖，即落入区间的测试点所占比例。以校准残差的经验分位构造区间的思路参照文献\cite{vovk2005conformal}；由于谱点连续相关，本文按留出点实计覆盖，检验对象是反射率预测，而非几何厚度。
% src:求解/问题2/结果/区间覆盖.json:测试点数;构造方法
% src:求解/问题3/结果/回炉轮3候选/仲裁闭环/20260910_131501_81623_建模公式/厚度结果.json:预测区间.4.经验覆盖率;预测区间.4.平均区间宽度_比例;预测区间.5.经验覆盖率;预测区间.5.平均区间宽度_比例;预测区间.6.经验覆盖率;预测区间.6.平均区间宽度_比例;预测区间.7.经验覆盖率;预测区间.7.平均区间宽度_比例;预测区间.20.经验覆盖率;预测区间.20.平均区间宽度_比例;预测区间.21.经验覆盖率;预测区间.21.平均区间宽度_比例;预测区间.22.经验覆盖率;预测区间.22.平均区间宽度_比例;预测区间.23.经验覆盖率;预测区间.23.平均区间宽度_比例;预测区间.36.经验覆盖率;预测区间.36.平均区间宽度_比例;预测区间.37.经验覆盖率;预测区间.37.平均区间宽度_比例;预测区间.38.经验覆盖率;预测区间.38.平均区间宽度_比例;预测区间.39.经验覆盖率;预测区间.39.平均区间宽度_比例;预测区间.52.经验覆盖率;预测区间.52.平均区间宽度_比例;预测区间.53.经验覆盖率;预测区间.53.平均区间宽度_比例;预测区间.54.经验覆盖率;预测区间.54.平均区间宽度_比例;预测区间.55.经验覆盖率;预测区间.55.平均区间宽度_比例;预测区间.4.名义覆盖率
% src:求解/问题3/结果/回炉轮3候选/仲裁闭环/20260910_131501_81623_建模公式/厚度结果.json:逐块评分.4.测试点数;逐块评分.5.测试点数;逐块评分.6.测试点数;逐块评分.7.测试点数;逐块评分.20.测试点数;逐块评分.21.测试点数;逐块评分.22.测试点数;逐块评分.23.测试点数;逐块评分.36.测试点数;逐块评分.37.测试点数;逐块评分.38.测试点数;逐块评分.39.测试点数;逐块评分.52.测试点数;逐块评分.53.测试点数;逐块评分.54.测试点数;逐块评分.55.测试点数

问题一的$19.54\,\mu\mathrm{m}$是周期等效量$T=5000/\Delta\sigma$，包含色散和界面相位变化的贡献，与后两问由光学模型求得的拟合厚度分列。各问的扰动比较记录厚度随光学条件的变化，文献提供通用方法，本题观测与模型计算支持样片结论。附录按问题列出计算材料、数据口径与完整比较，便于复核公式如何落实到观测与结果。
% src:求解/问题1/结果/模型验证.json:真实附件诊断.表观光学厚度乘积_dq_微米
% src:交接/结果声明_问题1.json:置信.等级;置信.理由
% src:交接/结果声明_问题2.json:置信.等级;置信.理由
% src:交接/结果声明_问题3.json:置信.等级;置信.理由

\begingroup
\renewcommand{\addcontentsline}[3]{}
\patchcmd{\thebibliography}{\section*{\refname}}{}{}{}
\begin{thebibliography}{9}
\bibitem{vovk2005conformal} Vovk V, Gammerman A, Shafer G. Algorithmic Learning in a Random World[M]. New York: Springer, 2005.
\end{thebibliography}
\endgroup
